\chapter{Cross-Language Validation}
\label{ch:validation}

\section{Introduction}

Implementation correctness is foundational to causal inference. A methodologically sound estimator produces misleading results if implemented incorrectly. This chapter describes the six-layer validation architecture used throughout this book.

\begin{definition}[Cross-Language Validation]
Independent implementations in different languages (Python, Julia) should produce identical results for identical inputs:
\begin{equation}
\left| \hat{\tau}_{\text{Python}} - \hat{\tau}_{\text{Julia}} \right| < 10^{-10}
\end{equation}
Agreement validates algorithmic correctness; divergence signals implementation errors.
\end{definition}

This chapter develops:
\begin{enumerate}
\item \textbf{Six-Layer Architecture}: Comprehensive validation framework
\item \textbf{Cross-Language Testing}: Python $\leftrightarrow$ Julia parity
\item \textbf{Monte Carlo Validation}: Statistical properties verification
\item \textbf{Adversarial Testing}: Edge case robustness
\end{enumerate}


\section{Six-Layer Validation Architecture}

\subsection{Architecture Overview}

\begin{table}[h]
\centering
\caption{Validation Layers}
\begin{tabular}{clll}
\toprule
Layer & Name & Purpose & Implementation \\
\midrule
1 & Known-Answer & Hand-calculated expected values & \texttt{tests/test\_*/*.py} \\
2 & Adversarial & Edge cases, boundary conditions & \texttt{tests/validation/adversarial/} \\
3 & Monte Carlo & Statistical properties (5K+ runs) & \texttt{tests/validation/monte\_carlo/} \\
4 & Cross-Language & Python $\leftrightarrow$ Julia parity & \texttt{tests/validation/cross\_language/} \\
5 & R Triangulation & External reference (future) & --- \\
6 & Golden Reference & Frozen reference results & \texttt{tests/golden\_results/} \\
\bottomrule
\end{tabular}
\end{table}

\subsection{Layer 1: Known-Answer Tests}

For simple cases with hand-calculable results:

\begin{minted}[fontsize=\small]{python}
def test_simple_ate_known_answer():
    """Four-unit case with known answer."""
    outcomes = np.array([7.0, 5.0, 3.0, 1.0])
    treatment = np.array([1, 1, 0, 0])

    # Hand calculation:
    # Y_1 = (7+5)/2 = 6.0
    # Y_0 = (3+1)/2 = 2.0
    # ATE = 6.0 - 2.0 = 4.0
    result = simple_ate(outcomes, treatment)

    assert result["estimate"] == 4.0, "ATE should be exactly 4.0"
\end{minted}

\subsection{Layer 2: Adversarial Tests}

Tests that probe edge cases and failure modes:

\begin{minted}[fontsize=\small]{python}
class TestSimpleATEAdversarial:
    """Edge cases that should fail gracefully."""

    def test_single_treated(self):
        """Only one treated unit (n_1 = 1)."""
        outcomes = np.array([10.0, 2.0, 3.0])
        treatment = np.array([1, 0, 0])

        # Should handle gracefully (large SE, wide CI)
        result = simple_ate(outcomes, treatment)
        assert result["se"] > 0

    def test_all_treated(self):
        """All units treated (no control group)."""
        outcomes = np.array([5.0, 6.0, 7.0])
        treatment = np.array([1, 1, 1])

        # Should raise informative error
        with pytest.raises(ValueError, match="control"):
            simple_ate(outcomes, treatment)

    def test_extreme_outlier(self):
        """Outlier 1000x larger than typical values."""
        outcomes = np.array([1.0, 2.0, 1e6, 3.0])
        treatment = np.array([1, 1, 0, 0])

        # Should compute without numerical issues
        result = simple_ate(outcomes, treatment)
        assert np.isfinite(result["estimate"])
\end{minted}

\subsection{Layer 3: Monte Carlo Validation}

Statistical validation over thousands of simulations:

\begin{minted}[fontsize=\small]{python}
@pytest.mark.monte_carlo
def test_simple_ate_unbiased():
    """Monte Carlo validation: 5000 runs, bias < 0.05, coverage 94-96%."""
    n_runs = 5000
    true_ate = 2.0

    estimates = []
    ci_covers = []

    for seed in range(n_runs):
        Y, T = generate_rct(n=100, true_ate=true_ate, seed=seed)
        result = simple_ate(Y, T)

        estimates.append(result["estimate"])
        ci_covers.append(result["ci_lower"] < true_ate < result["ci_upper"])

    # Bias check
    bias = np.mean(estimates) - true_ate
    assert abs(bias) < 0.05, f"Bias {bias:.4f} exceeds threshold"

    # Coverage check
    coverage = np.mean(ci_covers)
    assert 0.94 < coverage < 0.96, f"Coverage {coverage:.2%} outside range"
\end{minted}

\begin{table}[h]
\centering
\caption{Monte Carlo Validation Standards}
\begin{tabular}{lccc}
\toprule
Method Type & Bias Target & Coverage Target & SE Accuracy \\
\midrule
RCT (unconfounded) & $< 0.05$ & 94--96\% & $< 10\%$ \\
Observational & $< 0.10$ & 93--97\% & $< 15\%$ \\
PSM & $< 0.30$ & $\geq 95\%$ & $< 150\%$ \\
\bottomrule
\end{tabular}
\end{table}


\section{Cross-Language Validation}

\subsection{Design Principles}

1. \textbf{Independent implementations}: Python uses NumPy/SciPy; Julia uses native arrays
2. \textbf{Shared data}: Same random seeds and test fixtures
3. \textbf{Strict tolerance}: Agreement to 10 decimal places (rtol $< 10^{-10}$)
4. \textbf{Comprehensive coverage}: All estimators, not just simple cases

\subsection{Julia Interface Architecture}

\begin{minted}[fontsize=\small]{python}
# tests/validation/cross_language/julia_interface.py

from juliacall import Main as jl

# Initialize Julia project
JULIA_PROJECT_PATH = "julia/"
jl.seval(f'import Pkg; Pkg.activate("{JULIA_PROJECT_PATH}")')
jl.seval("using CausalEstimators")

def julia_simple_ate(outcomes, treatment, alpha=0.05):
    """Call Julia SimpleATE via juliacall."""
    # Convert to Julia types
    jl_outcomes = jl.collect(outcomes.astype(np.float64))
    jl_treatment = jl.collect(treatment.astype(bool))

    # Create problem and solve
    problem = jl.RCTProblem(jl_outcomes, jl_treatment)
    solution = jl.solve(problem, jl.SimpleATE())

    return {
        "estimate": float(solution.estimate),
        "se": float(solution.se),
        "ci_lower": float(solution.ci_lower),
        "ci_upper": float(solution.ci_upper),
    }
\end{minted}

\subsection{Cross-Language Test Pattern}

\begin{minted}[fontsize=\small]{python}
# tests/validation/cross_language/test_python_julia_simple_ate.py

import pytest
from tests.validation.cross_language.julia_interface import (
    is_julia_available,
    julia_simple_ate,
)

pytestmark = pytest.mark.skipif(
    not is_julia_available(),
    reason="Julia not available for cross-validation"
)

class TestPythonJuliaSimpleATE:
    """Cross-validate Python simple_ate against Julia SimpleATE."""

    def test_basic_case(self):
        """Basic 4-unit case."""
        outcomes = np.array([7.0, 5.0, 3.0, 1.0])
        treatment = np.array([1, 1, 0, 0])

        python_result = simple_ate(outcomes, treatment)
        julia_result = julia_simple_ate(outcomes, treatment)

        # Validate to 10 decimal places
        assert np.isclose(
            python_result["estimate"],
            julia_result["estimate"],
            rtol=1e-10
        ), f"Estimate mismatch: Python={python_result['estimate']}, Julia={julia_result['estimate']}"

        assert np.isclose(
            python_result["se"],
            julia_result["se"],
            rtol=1e-10
        ), f"SE mismatch: Python={python_result['se']}, Julia={julia_result['se']}"

    def test_large_sample(self):
        """Large sample (n=1000) stress test."""
        np.random.seed(42)
        n = 1000
        treatment = np.random.binomial(1, 0.5, n)
        outcomes = 2.0 * treatment + np.random.normal(0, 1, n)

        python_result = simple_ate(outcomes, treatment)
        julia_result = julia_simple_ate(outcomes, treatment)

        assert np.isclose(python_result["estimate"], julia_result["estimate"], rtol=1e-10)
        assert np.isclose(python_result["se"], julia_result["se"], rtol=1e-10)
\end{minted}

\subsection{Method Coverage}

\begin{table}[h]
\centering
\caption{Cross-Language Test Coverage}
\begin{tabular}{lll}
\toprule
Method Family & Python Module & Julia Module \\
\midrule
RCT & \texttt{rct.estimators} & \texttt{RCT.jl} \\
DiD & \texttt{did.*} & \texttt{did/*.jl} \\
IV & \texttt{iv.*} & \texttt{iv/*.jl} \\
RDD & \texttt{rdd.*} & \texttt{rdd/*.jl} \\
SCM & \texttt{scm.*} & \texttt{scm/*.jl} \\
CATE & \texttt{cate.*} & \texttt{cate/*.jl} \\
PSM & \texttt{psm.*} & \texttt{psm/*.jl} \\
Sensitivity & \texttt{sensitivity.*} & \texttt{sensitivity/*.jl} \\
QTE & \texttt{qte.*} & \texttt{qte/*.jl} \\
MTE & \texttt{mte.*} & \texttt{mte/*.jl} \\
Bounds & \texttt{bounds.*} & \texttt{bounds/*.jl} \\
Mediation & \texttt{mediation.*} & \texttt{mediation/*.jl} \\
\bottomrule
\end{tabular}
\end{table}


\section{Data Generating Processes}

\subsection{DGP Architecture}

Centralized DGP generators ensure identical data across tests:

\begin{minted}[fontsize=\small]{python}
# tests/validation/monte_carlo/dgp_generators.py

def dgp_simple_rct(n=100, true_ate=2.0, random_state=None):
    """Simple RCT with homoskedastic errors."""
    rng = np.random.default_rng(random_state)

    treatment = rng.binomial(1, 0.5, n)
    potential_Y0 = rng.normal(0, 1, n)
    potential_Y1 = potential_Y0 + true_ate
    outcomes = np.where(treatment == 1, potential_Y1, potential_Y0)

    return outcomes, treatment

def dgp_confounded_observational(n=500, true_ate=2.0, random_state=None):
    """Observational study with confounding."""
    rng = np.random.default_rng(random_state)

    # Confounder affects both treatment and outcome
    X = rng.normal(0, 1, n)
    propensity = 1 / (1 + np.exp(-0.5 * X))
    treatment = rng.binomial(1, propensity)

    # Outcome depends on X and treatment
    outcomes = 0.5 * X + true_ate * treatment + rng.normal(0, 1, n)

    return outcomes, treatment, X, propensity
\end{minted}

\subsection{Method-Specific DGPs}

\begin{minted}[fontsize=\small]{python}
# tests/validation/monte_carlo/dgp_did.py

def dgp_did_staggered(n_units=100, n_periods=10, true_ate=2.0, random_state=None):
    """Staggered DiD with heterogeneous timing."""
    rng = np.random.default_rng(random_state)

    # Staggered treatment adoption
    treatment_timing = rng.choice(range(3, n_periods), n_units)

    # Panel data with unit and time fixed effects
    unit_fe = rng.normal(0, 0.5, n_units)
    time_fe = rng.normal(0, 0.3, n_periods)

    # Build panel
    data = []
    for i in range(n_units):
        for t in range(n_periods):
            treated = 1 if t >= treatment_timing[i] else 0
            Y = unit_fe[i] + time_fe[t] + true_ate * treated + rng.normal(0, 0.5)
            data.append({"unit": i, "time": t, "treated": treated, "Y": Y})

    return pd.DataFrame(data)
\end{minted}


\section{Implementation Details}

\subsection{Running Validation Tests}

\begin{minted}[fontsize=\small]{bash}
# Full validation suite
pytest tests/validation/ -v --tb=short

# Layer-specific tests
pytest tests/validation/monte_carlo/ -v -m monte_carlo
pytest tests/validation/adversarial/ -v
pytest tests/validation/cross_language/ -v

# Skip slow Monte Carlo tests
pytest tests/ -m "not slow"

# Single method validation
pytest tests/validation/cross_language/test_python_julia_did.py -v
\end{minted}

\subsection{CI/CD Integration}

\begin{minted}[fontsize=\small]{yaml}
# .github/workflows/validation.yml
name: Validation Suite

on: [push, pull_request]

jobs:
  monte-carlo:
    runs-on: ubuntu-latest
    steps:
      - uses: actions/checkout@v4
      - name: Run Monte Carlo tests
        run: pytest tests/validation/monte_carlo/ -v --timeout=600

  cross-language:
    runs-on: ubuntu-latest
    steps:
      - uses: actions/checkout@v4
      - uses: julia-actions/setup-julia@v1
      - name: Install Julia deps
        run: julia --project=julia -e 'using Pkg; Pkg.instantiate()'
      - name: Run cross-language tests
        run: pytest tests/validation/cross_language/ -v
\end{minted}

\subsection{Debugging Validation Failures}

\begin{minted}[fontsize=\small]{python}
def debug_cross_language_failure(method_name, python_result, julia_result):
    """Diagnostic output for validation failures."""
    print(f"\n=== Cross-Language Failure: {method_name} ===")

    for key in python_result:
        py_val = python_result[key]
        jl_val = julia_result.get(key, "MISSING")

        if isinstance(py_val, (int, float)) and isinstance(jl_val, (int, float)):
            diff = abs(py_val - jl_val)
            rel_diff = diff / (abs(py_val) + 1e-10)
            status = "OK" if rel_diff < 1e-10 else "MISMATCH"
            print(f"  {key}: Python={py_val:.12e}, Julia={jl_val:.12e}, "
                  f"diff={diff:.2e} ({status})")
        else:
            print(f"  {key}: Python={py_val}, Julia={jl_val}")
\end{minted}


\section{Monte Carlo Standards}

\subsection{Validation Metrics}

\begin{definition}[Monte Carlo Validation Criteria]
For $R$ simulation runs with true effect $\tau$:
\begin{align}
\text{Bias} &= \frac{1}{R} \sum_{r=1}^R \hat{\tau}_r - \tau \\
\text{Coverage} &= \frac{1}{R} \sum_{r=1}^R \mathbf{1}[\hat{\tau}_r^L < \tau < \hat{\tau}_r^U] \\
\text{SE Accuracy} &= \frac{\left| \overline{\text{SE}} - \text{SD}(\hat{\tau}) \right|}{\text{SD}(\hat{\tau})}
\end{align}
\end{definition}

\subsection{Validation Helper}

\begin{minted}[fontsize=\small]{python}
# tests/validation/utils.py

def validate_monte_carlo_results(
    estimates,
    standard_errors,
    ci_lowers,
    ci_uppers,
    true_value,
    bias_threshold=0.05,
    coverage_range=(0.94, 0.96),
    se_accuracy_threshold=0.10,
):
    """Validate Monte Carlo simulation results."""
    estimates = np.array(estimates)
    standard_errors = np.array(standard_errors)
    ci_lowers = np.array(ci_lowers)
    ci_uppers = np.array(ci_uppers)

    # Bias
    bias = np.mean(estimates) - true_value
    bias_ok = abs(bias) < bias_threshold

    # Coverage
    coverage = np.mean((ci_lowers < true_value) & (true_value < ci_uppers))
    coverage_ok = coverage_range[0] < coverage < coverage_range[1]

    # SE accuracy
    empirical_se = np.std(estimates)
    mean_reported_se = np.mean(standard_errors)
    se_accuracy = abs(mean_reported_se - empirical_se) / empirical_se
    se_accuracy_ok = se_accuracy < se_accuracy_threshold

    return {
        "bias": bias,
        "bias_ok": bias_ok,
        "coverage": coverage,
        "coverage_ok": coverage_ok,
        "se_accuracy": se_accuracy,
        "se_accuracy_ok": se_accuracy_ok,
        "all_pass": bias_ok and coverage_ok and se_accuracy_ok,
    }
\end{minted}


\section{Golden Reference Tests}

\subsection{Frozen Reference Values}

Golden references provide regression testing against known-correct outputs:

\begin{minted}[fontsize=\small]{python}
# tests/golden_results/simple_ate.json
{
    "test_balanced_rct": {
        "input": {
            "outcomes": [7.0, 5.0, 3.0, 1.0],
            "treatment": [1, 1, 0, 0]
        },
        "expected": {
            "estimate": 4.0,
            "se": 1.4142135623730951,
            "ci_lower": -0.5319734684099524,
            "ci_upper": 8.531973468409952
        }
    }
}
\end{minted}

\begin{minted}[fontsize=\small]{python}
def test_against_golden():
    """Test against frozen reference values."""
    import json

    with open("tests/golden_results/simple_ate.json") as f:
        golden = json.load(f)

    for test_name, data in golden.items():
        result = simple_ate(
            np.array(data["input"]["outcomes"]),
            np.array(data["input"]["treatment"])
        )

        for key, expected in data["expected"].items():
            assert np.isclose(result[key], expected, rtol=1e-10), \
                f"{test_name}.{key}: got {result[key]}, expected {expected}"
\end{minted}


\section{Discussion}

\subsection{Key Takeaways}

\begin{enumerate}
\item \textbf{Six layers}: Comprehensive validation from unit tests to cross-language parity
\item \textbf{Cross-language}: Python and Julia must agree to 10 decimal places
\item \textbf{Monte Carlo}: 5,000+ runs verify statistical properties (bias, coverage, SE)
\item \textbf{Adversarial}: Edge cases test robustness and error handling
\item \textbf{Golden references}: Regression testing against frozen correct outputs
\end{enumerate}

\subsection{Practical Recommendations}

\begin{enumerate}
\item \textbf{Run full suite}: Before any release or major change
\item \textbf{Check Monte Carlo}: When modifying estimator internals
\item \textbf{Cross-language first}: When adding new estimators
\item \textbf{Update golden}: When correct behavior changes intentionally
\item \textbf{Document xfails}: Known limitations with explicit reasons
\end{enumerate}

\subsection{Validation Status by Method}

\begin{table}[h]
\centering
\caption{Current Validation Status}
\begin{tabular}{lccccc}
\toprule
Method & Known-Answer & Adversarial & Monte Carlo & Cross-Language & Status \\
\midrule
RCT & \checkmark & \checkmark & \checkmark & \checkmark & VERIFIED \\
DiD & \checkmark & \checkmark & \checkmark & \checkmark & VERIFIED \\
IV & \checkmark & \checkmark & \checkmark & \checkmark & VERIFIED \\
RDD & \checkmark & \checkmark & \checkmark & \checkmark & VERIFIED \\
SCM & \checkmark & \checkmark & \checkmark & \checkmark & VERIFIED \\
CATE & \checkmark & \checkmark & \checkmark & \checkmark & VERIFIED \\
PSM & \checkmark & \checkmark & \checkmark & \checkmark & VERIFIED \\
QTE & \checkmark & \checkmark & \checkmark & \checkmark & VERIFIED \\
MTE & \checkmark & \checkmark & \checkmark & \checkmark & VERIFIED \\
Mediation & \checkmark & \checkmark & \checkmark & \checkmark & VERIFIED \\
Selection & \checkmark & \checkmark & \checkmark & \checkmark & VERIFIED \\
DTR & \checkmark & \checkmark & \checkmark & \checkmark & VERIFIED \\
Principal Strat & \checkmark & \checkmark & \checkmark & \checkmark & VERIFIED \\
Bayesian & \checkmark & \checkmark & \checkmark & --- & PARTIAL \\
\bottomrule
\end{tabular}
\end{table}

\subsection{Limitations}

\begin{itemize}
\item \textbf{R triangulation}: Not yet implemented (planned Layer 5)
\item \textbf{Julia dependency}: Cross-language tests require juliacall installation
\item \textbf{Monte Carlo time}: Full suite takes $\sim$30 minutes
\item \textbf{Coverage gaps}: Some edge cases may not be tested
\end{itemize}

\subsection{Connection to Other Chapters}

\begin{itemize}
\item \textbf{All chapters}: Validation ensures correctness of presented implementations
\item \textbf{Appendix A}: Detailed DGP specifications for Monte Carlo
\item \textbf{Appendix B}: Complete test coverage matrix
\end{itemize}

